% PIRTM-ACE×PETC: Prime-Indexed Recursive Tensor Model with ACE Safety and PETC Lawfulness
\documentclass[11pt]{article}

% --- Packages ---
\usepackage[a4paper,margin=1in]{geometry}
\usepackage{amsmath,amsthm,amssymb,amsfonts,mathtools}
\usepackage{bm}
\usepackage{microtype}
\usepackage{enumitem}
\usepackage{hyperref}
\usepackage{cleveref}
\usepackage{algorithm}
\usepackage{algpseudocode}
\usepackage{listings}
\usepackage{xcolor}
\usepackage{graphicx}

% --- Hyperref setup ---
\hypersetup{
  colorlinks=true,
  linkcolor=blue,
  citecolor=blue,
  urlcolor=blue
}

% --- Listings (Python) ---
\definecolor{codebg}{RGB}{247,247,247}
\definecolor{codestring}{RGB}{163,21,21}
\definecolor{codecomment}{RGB}{0,128,0}
\definecolor{codekeyword}{RGB}{0,0,180}
\lstdefinestyle{py}{
  language=Python,
  basicstyle=\ttfamily\small,
  backgroundcolor=\color{codebg},
  keywordstyle=\color{codekeyword}\bfseries,
  stringstyle=\color{codestring},
  commentstyle=\color{codecomment}\itshape,
  showstringspaces=false,
  frame=single,
  framerule=0.4pt,
  rulecolor=\color{black},
  columns=fullflexible,
  keepspaces=true,
  tabsize=2,
  breaklines=true,
  numbers=left,
  numbersep=6pt,
  numberstyle=\tiny\color{gray}
}

% --- Theorem Environments ---
\theoremstyle{plain}
\newtheorem{theorem}{Theorem}
\newtheorem{proposition}[theorem]{Proposition}
\newtheorem{lemma}[theorem]{Lemma}
\newtheorem{corollary}[theorem]{Corollary}
\theoremstyle{definition}
\newtheorem{definition}[theorem]{Definition}
\theoremstyle{remark}
\newtheorem{remark}[theorem]{Remark}

% --- Macros ---
\newcommand{\R}{\mathbb{R}}
\newcommand{\N}{\mathbb{N}}
\newcommand{\norm}[1]{\left\lVert #1 \right\rVert}
\newcommand{\abs}[1]{\left\lvert #1 \right\rvert}
\newcommand{\braces}[1]{\left\{ #1 \right\}}
\newcommand{\angles}[1]{\left\langle #1 \right\rangle}
\newcommand{\op}{\mathrm{op}}
\newcommand{\Id}{\mathrm{Id}}
\newcommand{\Kop}{\mathcal{K}}

\begin{document}


\title{PIRTM --- Core Technical Disclosure \\ Defensive Publication / Prior Art}
\author {Ryan Van Gelder (PETC) \& Tyler Van Osdol (ACE) \\ \\ Citizen Gardens}
\date{\today}
\maketitle

\begin{abstract}
We present a complete, engineering-ready formulation of the Prime-Indexed Recursive Tensor Model (PIRTM) integrated with the Arithmetic Control Engine (ACE) for safety certificates and the Prime-Encoded Tensor Calculus (PETC) for lawful prime-weighted structure. The core dynamics are affine-linear with a prime-decomposed gain operator. We prove necessary and sufficient convergence conditions (spectral radius $<1$), provide sufficient ACE--PETC budget tests, give robustness and time-varying extensions, and supply drop-in Python implementations including exact weighted-$\ell_1$ projection, monitoring utilities, fixed-point estimators with certified tails, and adaptive margin control. Throughout we adhere to the notational convention that the gain acts by operator application $K\,T$ (\emph{not} tensor product), reserving $\otimes$ for true space tensor products.
\end{abstract}
\newpage
\tableofcontents

\section{Introduction}
PIRTM models the evolution of a tensor-valued state $T_t$ under a prime-weighted linear operator with an affine forcing term. To guarantee safe and predictable behavior, we combine:
\begin{itemize}[leftmargin=*]
  \item \textbf{PETC}: a structural layer that encodes prime indices, exponents and operator bounds $b_p$, producing a lawful ledger for $\{(p,\alpha,b_p)\}$ and raw weights;
  \item \textbf{ACE}: a certification layer that enforces a contraction budget via exact weighted-$\ell_1$ projection and yields online certificates (contraction gap, rate upper bounds).
\end{itemize}
This report consolidates the mathematical foundation (convergence, robustness) and the production-grade implementation (projection, monitoring and solvers).

\section{Notation and Setting}
Let $(H,\norm{\cdot})$ be a Banach space suitable for the tensor $T_t\in H$. For a list of primes $P_N=\{p_1,\dots,p_N\}$ and bounded linear maps $B_p:H\to H$ with operator bounds $\norm{B_p}\le b_p$, define prime weights
\[
  w_p := \Lambda_p\,p^{\alpha}, \qquad \alpha\in\R, \ \Lambda_p\in\R.
\]
The gain operator is the (finite) sum
\[
  K := \sum_{p\in P_N} w_p\,B_p, \qquad \text{acting by application } (K\,X)=\sum_{p} w_p\,B_p(X).
\]
We reserve $\otimes$ for explicit tensor products of spaces; we \emph{do not} write $K\otimes T$ for the state update.

\section{Core Recurrence: Scalar and Operator Forms}
\subsection{Scalar/Mode-wise Form}
For each mode $(m,n)$, the state obeys
\begin{equation}
  T_{t+1}(m,n) = k_m\,T_t(m,n) + F(m,n), \qquad k_m := \sum_{p\in P_N} \Lambda_m\, p^{\alpha}.
  \label{eq:scalar}
\end{equation}
This is a specialization of the operator model when each $B_p$ acts as the identity on the $(m,n)$-mode.

\subsection{Operator Form}
For the full tensor state $T_t\in H$ we have
\begin{equation}
  T_{t+1} = F + K\,T_t, \qquad K = \sum_{p\in P_N} w_p\,B_p, \quad w_p=\Lambda_p\,p^{\alpha}.
  \label{eq:operator}
\end{equation}

\section{Convergence Theory}
\begin{theorem}[Fixed point and convergence]
Consider \eqref{eq:operator}. The following are equivalent:
\begin{enumerate}[label=(\roman*),leftmargin=*]
  \item The spectral radius satisfies $\rho(K)<1$.
  \item The operator $(I-K)$ is invertible with Neumann series $(I-K)^{-1}=\sum_{j\ge 0} K^j$ converging in operator norm.
  \item For every $T_0\in H$, the sequence $T_{t+1}=F+K\,T_t$ converges to the unique fixed point $T_\infty=(I-K)^{-1}F$ with linear rate bounded by $\norm{K}^t$.
\end{enumerate}
\end{theorem}
\begin{proof}
$(i)\Rightarrow(ii)$: If $\rho(K)<1$, then there exists a submultiplicative operator norm with $\norm{K}<1$, which yields the Neumann series. $(ii)\Rightarrow(iii)$: Summation of $K^jF$ and the Banach fixed-point theorem give uniqueness and linear convergence. $(iii)\Rightarrow(i)$: If iterates converge for all $T_0$ then $1$ is not in the spectrum of $K$ and $\rho(K)<1$.
\end{proof}

\begin{corollary}[Scalar specialization]
If each $B_p=\Id$, then $K=k\,\Id$ with $k=\sum_p w_p$. Convergence holds iff $|k|<1$, and the fixed point and rate are
\[
  T_\infty = \frac{F}{1-k}, \qquad |T_t-T_\infty| \le |k|^t\,|T_0-T_\infty|.
\]
\end{corollary}

\subsection{ACE--PETC Sufficient Test}
A practical sufficient condition enforced by ACE with PETC bounds is
\begin{equation}
  \sum_{p\in P_N} b_p\,|w_p|\; <\; 1-\varepsilon, \qquad (\varepsilon>0),
  \label{eq:ace-budget}
\end{equation}
which implies $\norm{K}\le \sum_p b_p|w_p| < 1$ and hence $\rho(K)<1$.

\subsection{Infinite-prime extension}
If $P_N\uparrow P$ and $\sum_{p\in P} b_p\,|\Lambda_p|\,p^{\alpha} < 1$ (e.g., $\alpha<-1$ with bounded $\Lambda_p$ and $b_p$), then the series for $K$ converges absolutely in operator norm and all results above hold.

\section{Robustness and Time-Varying Systems}
\begin{proposition}[Sensitivity of the fixed point]
Let $\norm{K}\le 1-\delta$ with $\delta\in(0,1)$. For perturbations $\Delta K$ and $\Delta F$ with $\norm{\Delta K}<\delta$, the fixed point shifts by
\[
  \Delta T_\infty = (I-K)^{-1}\,\Delta F + (I-K)^{-1}\,\Delta K\,(I-K)^{-1}F + \mathcal{O}(\norm{\Delta K}^2),
\]
so that $\norm{\Delta T_\infty}\le \delta^{-1}\norm{\Delta F} + \delta^{-2}\norm{\Delta K}\,\norm{F}+\ldots$.
\end{proposition}

\begin{theorem}[Uniform contraction for time-varying gains]
Suppose $T_{t+1}=F_t+K_t T_t$ with $\sup_t\norm{K_t}\le 1-\delta$ and bounded $F_t$. Then
\[
  \norm{T_t}\le (1-\delta)^t\,\norm{T_0} + \frac{1}{\delta}\,\sup_{j\le t}\norm{F_j},
\]
so the system is BIBO-stable and tracks slowly varying fixed points with error $\mathcal{O}(\delta^{-1})$.
\end{theorem}

\section{ACE$\times$PETC Integration}
\paragraph{PETC (lawfulness).} PETC supplies the prime-signature ledger $\{(p,\alpha,b_p)\}$ and raw weights $w_p=\Lambda_p p^{\alpha}$. It also specifies admissible decay profiles in $\Lambda_p$ when $\alpha\ge -1$ so that \eqref{eq:ace-budget} still holds.

\paragraph{ACE (safety).} ACE enforces the weighted-$\ell_1$ budget \eqref{eq:ace-budget} via a 1-Lipschitz projection (soft-thresholding by a common multiplier), preserving contraction margins. Online metrics: the contraction upper bound $\widehat{\norm{K}}=\sum_p b_p|w_p|$, the gap $\mathrm{gapLB}=1-\widehat{\norm{K}}$, and a rate upper bound $\mathrm{SlopeUB}=\widehat{\norm{K}}$.

\section{Algorithms and Reference Implementations}
\subsection{Exact weighted-$\ell_1$ projection (ACE)}
\begin{algorithm}[H]
\caption{Weighted-$\ell_1$ projection by bisection}
\begin{algorithmic}[1]
\Require raw weights $w$, PETC bounds $b$, budget $\tau\in(0,1)$
\If{$\sum_i b_i|w_i|\le \tau$} \State \Return $w$ \EndIf
\State Set $\lambda\in[0,\max_i |w_i|/\max(b_i,\epsilon)]$
\While{budget error $>\,$tol and iters $<\,$max}
  \State $w^{\mathrm{proj}}_i \leftarrow \mathrm{sign}(w_i)\,\max(0, |w_i|-\lambda b_i)$
  \State $\mathrm{budget}\leftarrow \sum_i b_i |w^{\mathrm{proj}}_i|$
  \State Adjust $\lambda$ by bisection to meet $\mathrm{budget}=\tau$
\EndWhile
\State \Return $w^{\mathrm{proj}}$
\end{algorithmic}
\end{algorithm}

\paragraph{Python.}
\begin{lstlisting}[style=py,caption={Exact weighted-$\ell_1$ projection (ACE budget)}]
from typing import List
import numpy as np

def project_weighted_l1(w: List[float], b: List[float], tau: float, 
                        tol: float = 1e-9, max_iters: int = 80) -> List[float]:
    """
    Projects weights to exactly meet ACE budget via bisection.
    Returns projected w satisfying sum(b_i * |w_i|) ≈ tau.
    """
    current_budget = sum(b_i * abs(w_i) for w_i, b_i in zip(w, b))
    if current_budget <= tau:
        return w.copy()

    lo, hi = 0.0, max(abs(w_i) / max(b_i, 1e-15) for w_i, b_i in zip(w, b))
    w_proj = w.copy()
    for _ in range(max_iters):
        lam = 0.5 * (lo + hi)
        w_proj = [
            np.sign(w_i) * max(0.0, abs(w_i) - lam * b_i)
            for w_i, b_i in zip(w, b)
        ]
        budget = sum(b_i * abs(wi) for wi, b_i in zip(w_proj, b))
        if abs(budget - tau) < tol:
            break
        if budget > tau:
            lo = lam
        else:
            hi = lam
    return w_proj
\end{lstlisting}

\subsection{Unified ACE--PETC update (scalar and operator cases)}
\begin{lstlisting}[style=py,caption={One PIRTM step with ACE×PETC certification}]
from typing import List, Callable, Dict, Tuple
import numpy as np

Tensor = np.ndarray
LinearOp = Callable[[Tensor], Tensor]


def ace_petc_step(T: Tensor, F: Tensor, primes: List[int], alpha: float,
                  Lambda: List[float], B_ops: List[LinearOp] = None,
                  b_p: List[float] = None, tau: float = 0.9
                  ) -> Tuple[Tensor, Dict, Dict]:
    """
    One step of PIRTM recurrence with ACE×PETC certification.
    Returns: T_next, margins, debug
    """
    w_raw = [Lam_i * (p ** alpha) for p, Lam_i in zip(primes, Lambda)]

    if B_ops is None:
        # Scalar specialization
        b_scalar = b_p if b_p is not None else [1.0] * len(primes)
        w_proj = project_weighted_l1(w_raw, b_scalar, tau)
        k_eff = sum(w_proj)
        T_next = k_eff * T + F
        margins = {
            "||K||": abs(k_eff),
            "gapLB": 1.0 - abs(k_eff),
            "contraction_rate": abs(k_eff),
            "budget_used": sum(b_i * abs(w_i) for w_i, b_i in zip(w_proj, b_scalar)),
        }
        debug = {"w": w_proj, "k_eff": k_eff, "w_raw": w_raw}
        return T_next, margins, debug

    # Operator case
    assert b_p is not None and len(b_p) == len(primes)
    w_proj = project_weighted_l1(w_raw, b_p, tau)

    def K_operator(X: Tensor) -> Tensor:
        Y = np.zeros_like(X)
        for w_i, B_i in zip(w_proj, B_ops):
            Y = Y + w_i * B_i(X)
        return Y

    T_next = K_operator(T) + F
    K_norm_upper = sum(abs(w_i) * bp for w_i, bp in zip(w_proj, b_p))
    margins = {
        "||K||_ub": K_norm_upper,
        "gapLB": 1.0 - K_norm_upper,
        "contraction_rate": K_norm_upper,
        "budget_used": K_norm_upper,
    }
    debug = {"w": w_proj, "K_norm_ub": K_norm_upper, "K_operator": K_operator}
    return T_next, margins, debug
\end{lstlisting}

\subsection{Monitoring and safety status}
\begin{lstlisting}[style=py,caption={Real-time ACE margins monitor}]
class PIRTMMonitor:
    """Real-time monitoring of ACE×PETC safety margins"""
    def __init__(self, tau: float = 0.9, window: int = 100, norm_fn=None):
        self.tau = tau
        self.window = window
        self.norm_fn = norm_fn if norm_fn is not None else \
            (lambda x: float(np.linalg.norm(np.asarray(x))))
        self.history = {"K_norm": [], "gapLB": [], "T_diff": [], "sensitivity": []}

    def update(self, margins: Dict, T_prev: Tensor, T_current: Tensor):
        K_norm = margins.get("||K||", margins.get("||K||_ub", 0.0))
        gap = margins.get("gapLB", 0.0)
        self.history["K_norm"].append(K_norm)
        self.history["gapLB"].append(gap)
        self.history["T_diff"].append(self.norm_fn(np.asarray(T_current) - np.asarray(T_prev)))
        self.history["sensitivity"].append(1.0 / max(gap, 1e-9))
        for key in self.history:
            if len(self.history[key]) > self.window:
                self.history[key].pop(0)

    def safety_status(self) -> Dict:
        if not self.history["K_norm"]:
            return {"status": "NO_DATA"}
        recent_K = float(np.mean(self.history["K_norm"][ -5: ]))
        recent_gap = float(np.mean(self.history["gapLB"][ -5: ]))
        if recent_K > self.tau:
            status = "VIOLATION"
        elif recent_gap < 0.05:
            status = "CRITICAL"
        elif recent_gap < 0.10:
            status = "WARNING"
        else:
            status = "SAFE"
        return {
            "status": status,
            "avg_contraction": recent_K,
            "avg_gap": recent_gap,
            "sensitivity": 1.0 / max(recent_gap, 1e-9),
        }
\end{lstlisting}

\subsection{Fixed-point estimator with certified tail}
\begin{lstlisting}[style=py,caption={(I - K)^{-1}F via Neumann series with tail certificate}]
def fixed_point_estimate(F: Tensor,
                         K_apply: LinearOp,
                         K_norm_ub: float,
                         tol: float = 1e-8,
                         max_terms: int = 10_000) -> Tensor:
    """Computes (I - K)^{-1} F via Neumann series; requires K_norm_ub < 1."""
    assert K_norm_ub < 1.0, "Neumann series requires ||K|| < 1."
    S = np.zeros_like(F)
    term = F.copy()
    tau = K_norm_ub
    for j in range(max_terms):
        S = S + term
        tail_bound = (tau ** (j + 1)) * float(np.linalg.norm(np.asarray(F))) / (1 - tau)
        if tail_bound <= tol:
            break
        term = K_apply(term)
    return S
\end{lstlisting}

\subsection{Adaptive budget to maintain a target gap}
\begin{lstlisting}[style=py,caption={Simple adaptive policy for τ to keep a safety gap}]
def adapt_tau(tau: float, gapLB: float, gap_target: float = 0.1, eta: float = 0.5) -> float:
    """If gapLB < target, move tau toward 1 - gap_target by factor eta."""
    if gapLB < gap_target:
        return (1 - eta) * tau + eta * (1.0 - gap_target)
    return tau
\end{lstlisting}

\subsection{Infinite-prime convergence check (corrected)}
\begin{lstlisting}[style=py,caption={Absolute convergence test for ∑ b_p|Λ_p|p^α}]
def check_infinite_prime_convergence(alpha: float, Lambda_bound: float, b_bound: float) -> bool:
    """Returns True iff α < -1 and bounds are finite and nonnegative."""
    ok_bounds = np.isfinite(Lambda_bound) and np.isfinite(b_bound) and (Lambda_bound >= 0) and (b_bound >= 0)
    return (alpha < -1) and ok_bounds
\end{lstlisting}

\subsection{Usage example (scalar case)}
\begin{lstlisting}[style=py,caption={Scalar PIRTM loop with monitoring and fixed-point estimate}]
primes = [2, 3, 5, 7, 11]
alpha = -1.5
Lambda = [0.1] * len(primes)
F = 1.0
T = 0.0

monitor = PIRTMMonitor(tau=0.85)
for step in range(100):
    T_next, margins, debug = ace_petc_step(T, F, primes, alpha, Lambda, tau=0.85)
    monitor.update(margins, T, T_next)
    T = T_next
    if step % 10 == 0:
        status = monitor.safety_status()
        print(f"Step {step}: {status['status']} (K={margins.get('||K||', 0):.3f}, gap={margins['gapLB']:.3f})")

# Fixed point
k_eff = debug['k_eff']
T_inf = F / (1 - k_eff)
print(f"Fixed point: {T_inf:.6f}")
\end{lstlisting}

\subsection{Design Guidelines and Checklists}
\begin{itemize}[leftmargin=*]
  \item \textbf{Choose $\alpha$}: start with $\alpha\in(-2,-1)$; adjust $\Lambda_p$ to keep $\sum b_p|w_p|\le \tau$.
  \item \textbf{Projection}: always use the exact weighted-$\ell_1$ projection (not uniform scaling).
  \item \textbf{Monitoring}: log $\widehat{\norm{K}}$, gap $1-\widehat{\norm{K}}$, and $\norm{T_{t+1}-T_t}$.
  \item \textbf{Safety margins}: begin with $\tau\le 0.9$; adapt using \texttt{adapt\_tau} only when margins erode.
  \item \textbf{Infinite primes}: ensure $\sum_p b_p|\Lambda_p|p^{\alpha}<1$ (e.g., $\alpha<-1$ with bounded $\Lambda_p$, $b_p$), noting this is sufficient but not necessary.
\end{itemize}

\subsection{Conclusions}
The PIRTM dynamics combined with ACE and PETC yield a contractive, well-conditioned, and monitorable system. The theory (spectral radius criterion, budget-based sufficiency) aligns exactly with the provided implementations (projection, monitors, fixed-point estimation), enabling reliable deployment.

\section{Lightweight Scalar Friendly}
We study a prime-indexed linear iteration for tensors,
\begin{equation*}
  T_{t+1}=F+k\,T_t, \quad k\in\R,
\end{equation*}
and show that convergence is guaranteed \emph{iff} $|k|<1$. ``Prime structure'' is treated as a \emph{dictionary} (representation/regularization) and as \emph{typing metadata} via the \textbf{Prime-Encoded Tensor Calculus (PETC; Ryan Van Gelder)} \cite{VanGelderPETC}, rather than as a physical law. Stability and safety are certified with a 1-Lipschitz projection from the \textbf{Arithmetic Control Engine (ACE; Tyler Van Osdol)} \cite{VanOsdolACE}. We replace universal constants with explicit budgets, move speculative physics to future work with a parameter-identification plan, and add a preregistered experimental protocol.
\end{abstract}

\subsection{Notation and Setup}
Let $(\Hhh,\|\cdot\|)$ be a real Banach space (e.g., $\ell_2$ coefficients under a fixed dictionary). Fix a finite prime set $P_N=\{p_1,\dots,p_N\}$. A \emph{prime dictionary} is a collection $\D=\{\phi_p\}_{p\in P_N}\subset\Hhh$ used for expansions and penalties; we make no basis/completeness claim. For external drive $F\in\Hhh$, define the scalar gain
\begin{equation}
  k\;=\;\sum_{p\in P_N}\Lambda_p\,p^{\alpha},\qquad \Lambda_p\in\R,\;\alpha\in\R.
  \label{eq:gain}
\end{equation}
Older drafts used a single $\Lambda$; here we allow per-prime $\Lambda_p$.

\subsection{Convergence and Safety}
\begin{theorem}[Convergence $\Leftrightarrow$ contraction]\label{thm:contraction}
Let $\mathcal T(X)=F+kX$ on $\Hhh$. If $|k|<1$, then for any $T_0$,
\begin{equation*}
  T_t \to T_\infty = \frac{1}{1-k}\,F,
\end{equation*}
with linear rate $|k|$.
\end{theorem}
\begin{proof}
We have $\|\mathcal T(X)-\mathcal T(Y)\|=|k|\,\|X-Y\|$. If $|k|<1$, $\mathcal T$ is a contraction; apply Banach's fixed-point theorem \cite{BauschkeCombettes2017}.
\end{proof}

\begin{lemma}[One sufficient route to $|k|<1$]\label{lem:sufficient}
If $\Lambda_p\ge 0$ and $\alpha<-1$, then
\begin{equation*}
  |k|=\sum_{p\in P_N}\Lambda_p p^{\alpha}
  \;\le\; \Big(\max_{p}\Lambda_p\Big)\sum_{i=1}^{N}p_i^{\alpha}.
\end{equation*}
Choosing $\max_p\Lambda_p$ small enough makes $|k|<1$. \emph{(Sufficient, not necessary.)}
\end{lemma}

\paragraph{ACE safety wrapper (spectral certification).} Let $S\subset\Hhh$ be a safety/regularization set (e.g., a weighted-$\ell_1$ ball on dictionary coefficients). Let $\Pi_S$ be the Euclidean projection onto $S$ (1-Lipschitz). The ACE-wrapped update
\begin{equation}
  T_{t+1} \,=\, \Pi_S\!\big(F+k\,T_t\big)
  \label{eq:ace}
\end{equation}
has global Lipschitz constant $\le |k|$; thus $|k|<1$ gives a certified contraction and feasibility $T_t\in S$ for all $t$.

\subsection{PETC Typing: Prime Structure as Metadata}
Expand $T=\sum_{p\in P_N} w_p\,\phi_p$. Each coefficient $w_p$ carries a PETC ``prime type'' tag \cite{VanGelderPETC}. We impose a \emph{prime-mix budget}
\begin{equation}
  \sum_{p\in P_N} b_p\,|w_p(T)| \;\le\; \tau,\qquad b_p>0,
  \label{eq:budget}
\end{equation}
which is enforced via the projection $\Pi_S$ (weighted soft-thresholding on $w_p$). Optional \emph{prime gating} $g_p(\alpha,\Lambda_p)=\Lambda_p p^{\alpha}$ modulates proposals but is bounded by the same budget; no ``universal constant'' is assumed.

\subsection{Algorithm (ACE$\times$PETC-compliant)}
\noindent\textbf{Inputs.} $F$, dictionary $\{\phi_p\}$, weights $\{\Lambda_p\}$, exponent $\alpha$, budget $(b_p,\tau)$.

\begin{enumerate}[leftmargin=2em,label=\arabic*.]
  \item \textbf{Gain synthesis.} Compute $k$ from \cref{eq:gain}. If $|k|\ge 1$, shrink $\Lambda_p$ and/or $P_N$ until $|k|<1$.
  \item \textbf{Projection set.} Define $S=\{T:\sum_p b_p|w_p(T)|\le \tau\}$. Use $\Pi_S$ (1-Lipschitz).
  \item \textbf{Iteration.} $T_{t+1}=\Pi_S(F+kT_t)$ as in \cref{eq:ace}.
  \item \textbf{Stopping.} Stop when $\|T_{t+1}-T_t\|\le\varepsilon$; output $T_\infty$.
\end{enumerate}
\noindent\emph{Properties.} Certified convergence (\cref{thm:contraction}), feasibility (ACE), and auditable prime composition (PETC).

\subsection{Practical Guidance}
\paragraph{Choosing $\alpha,\Lambda_p,P_N$.} Pick $P_N$ by model/compute budget; start with $\alpha\in(-2,-1)$ and uniform $\Lambda_p=\Lambda$; use \cref{lem:sufficient} to get a safe $\Lambda$ with margin $\varepsilon\in[0.1,0.3]$; tune $(b_p,\tau)$ and re-check $|k|<1$.
\paragraph{Diagnostics.} Track $|k|$, budget usage $\sum b_p|w_p|$, and prime-mix histograms per iteration. If convergence slows, reduce $|k|$, tighten $\tau$, or shrink $P_N$.

\subsection{Experimental Protocol (Preregistered)}
\textbf{Benchmarks.} MSD-PrimeMix (multi-scale disturbances with controlled prime-spread).\\
\textbf{Baselines.} LQR (if model known), projected gradient with fixed step, Adam; RL baseline (e.g., TD3/DQN as applicable).\\
\textbf{Metrics.} Task cost, settling time, safety-violation count, robustness vs. disturbance amplitude/sparsity; report mean$\pm$95\% CI over 10 seeds.\\
\textbf{Ablations.} With/without ACE; with/without PETC budget; vary $\alpha$ at fixed $|k|$; expand/shrink $P_N$.\\
\textbf{Reporting.} Publish code/configs; include certificate traces (GapLB, SlopeUB).

\subsection{Related Work}
Contractions and projection operators are classical \cite{BauschkeCombettes2017}; proximal methods and Lipschitz projections connect to the modern optimization toolkit \cite{ParikhBoyd2014}. Weighted $\ell_1$ projections and soft-thresholding date to wavelet shrinkage \cite{DonohoJohnstone1994} and LASSO \cite{Tibshirani1996}. Prime-indexed dictionaries fit within sparse coding and dictionary learning \cite{MallatZhang1993,OlshausenField1996,Mairal2010}. Lipschitz-controlled neural networks relate via Parseval regularization \cite{Cisse2017Parseval} and spectral normalization \cite{Miyato2018SN}. The ACE and PETC frameworks formalize safety via contraction certificates and structural typing, respectively \cite{VanOsdolACE,VanGelderPETC}.

\subsection{Limitations}
The linear core is intentional; nonlinear variants require fresh contraction proofs (local or global). Dictionary choice is application-dependent; performance hinges on whether $\{\phi_p\}$ captures salient structure.

\subsection{Future Work}
\textbf{Operator learning.} Learn diagonal-plus-low-rank operators in the prime dictionary with PETC regularizers; verify a contraction region.\\
\textbf{Physical identification.} Units-consistent maps from data to $(\Lambda_p,\alpha)$ with uncertainty quantification.\\
\textbf{Nonlinear plants.} ACE-style sector-bounded contractions and piecewise certificates.
\section{The Prime-Recursive Foundations of Mathematical Existence}
\label{sec:prime_recursive_foundations}

\medskip

A central question is whether diverse structures—algebraic, topological, or computational—can be generated by a single, prime-based, recursive mechanism. We formalize this as a meta-recursive \emph{prime-indexed} constructor that is both \emph{auditable} (PETC: unambiguous retrospection) and \emph{stable} (ACE: certified contraction).

\subsection{The Meta-Recursive Function (ACE\texorpdfstring{$\times$}{x}PETC Form)}
Let $P_N=\{p_1,\dots,p_N\}$ be the first $N$ primes. Let $H$ be a Banach space of (typed) tensors. For each prime $p\in P_N$ choose a bounded linear channel $B_p:H\!\to\!H$ with $\|B_p\|\le b_p$. Define prime weights
\[
w_p=\Lambda_p\,p^{\alpha},\quad \Lambda_p\in\mathbb{R}\ (\text{typically }\Lambda_p\ge0),\quad \alpha< -1\ \text{(sufficient decay)}.
\]
The \emph{gain operator} and the \emph{meta-recursion} are
\begin{equation}
\label{eq:gain_and_meta}
\mathcal{K}=\sum_{p\in P_N} w_p B_p,\qquad
T_{t+1}=\mathcal{M}(P_N;T_t):=F+\mathcal{K}\,T_t.
\end{equation}
This refines the informal sum $\sum_{p} \Lambda\,p^\alpha T_{p}^{(m,n)}$ into an operator form while keeping primes explicit as independent channels.

\begin{assumption}[ACE Contraction Budget]
\label{as:budget}
\textstyle \sum_{p\in P_N} |w_p|\,b_p \;<\;1 \quad\Longrightarrow\quad \|\mathcal{K}\|\le\sum_{p}|w_p|b_p<1.
\end{assumption}

\begin{theorem}[Convergence \& Fixed Point]
\label{thm:convergence}
Under Assumption~\ref{as:budget}, the affine map $T\mapsto F+\mathcal{K}T$ is a strict contraction on $(H,\|\cdot\|)$. Hence there exists a unique fixed point
\[
T_\infty=(I-\mathcal{K})^{-1}F,
\]
and for any $T_0$ one has linear convergence
$\|T_t-T_\infty\|\le \|\mathcal{K}\|^{\,t}\,\|T_0-T_\infty\|$.
In the scalar special case ($T\in\mathbb{R}$, $\mathcal{K}=k\in(-1,1)$) this reduces to $T_\infty=\frac{F}{1-k}$.
\end{theorem}

\begin{theorem}[Uniqueness \& Stability under Projection]
\label{thm:uniqueness}
Let $\Pi_S$ be a nonexpansive projection (e.g., a convex projection) and consider the projected recursion
\[
T_{t+1}=\Pi_S\!\big(F+\mathcal{K}T_t\big).
\]
If Assumption~\ref{as:budget} holds, then $T\mapsto \Pi_S(F+\mathcal{K}T)$ is also a contraction with constant $\le\|\mathcal{K}\|<1$; therefore it admits a unique fixed point in $S$ and the iteration converges linearly to it.
\end{theorem}

\begin{remark}[Safe Design Regime for Prime Weights]
A sufficient (not necessary) recipe: choose $\alpha< -1$, $\Lambda_p\ge0$, and cap $\Lambda_p\le\bar\Lambda$ so that $\sum_{p\in P_N}\Lambda_p p^{\alpha} b_p<1$. Increasing $N$ or tightening $\bar\Lambda$ preserves the margin.
\end{remark}

\subsection{PETC: Prime-Signature Ledger for Unambiguous Retrospection}
Let $\mathbb{P}$ be the set of primes. A \emph{prime signature} is a finitely supported vector $e=(e_p)_{p\in\mathbb{P}}\in\mathbb{Z}^{(\mathbb{P})}$. Define the multiplicity map
\[
M(e)=\prod_{p} p^{e_p}\in\mathbb{Q}_{>0}.
\]
Composition of interactions is $e+e'$ (multiset union); undo/consumption is $e-e'$; dualization is $-e$.
Unique factorization makes the history canonical: $e_p=v_p(M(e))$ records how many times the $p$-tagged event occurred.

\begin{definition}[Conservation Certificate]
For any typed operation with inputs $X_1,\dots,X_r$ and output $Y$, PETC conservation requires
\(
e(Y)=\sum_{i=1}^r e(X_i).
\)
A system is \emph{ledger-consistent} if the equality holds for all executed operations.
\end{definition}

\begin{proposition}[Lossless Retrospection]
If a system is ledger-consistent, then for every state the accumulated signature $e$ determines the exact multiset of prime-tagged interactions that occurred, via $e_p=v_p(M(e))$.
\end{proposition}

\paragraph{Complexity.}
Maintain $e$ as a sparse map (prime $\mapsto$ exponent). Updates/audits run in $O(|\mathrm{supp}(e)|)$—no integer factorization is needed.

\subsection{Interpretation as a Universal Constructor}
\label{subsec:universal_constructor}
\textbf{Recursive self-modification.} The assignment $T_{t+1}=\mathcal{M}(P_N;T_t)$ mirrors prior states via $\mathcal{K}$, with uniqueness/convergence guaranteed by Theorems~\ref{thm:convergence}–\ref{thm:uniqueness}.

\textbf{Structure emergence.} Choosing $F$ (drive) and $\{B_p\}$ (channels) yields a wide class of constructs: neural weights, quantum states, algebraic operators, or kernel process states.

\textbf{Universality (within $H$).} Iteration converges to $T_\infty=(I-\mathcal{K})^{-1}F$; by selecting $F$ and $\{B_p\}$ appropriately, one can encode a broad family of objects in tensor space.

\subsection{Examples of Emergent Structures}
\begin{enumerate}
  \item \textbf{Algebraic Systems.} Let $B_p$ be structured projectors/sparsity masks; $F$ encodes a group/semiring operation tensor. Contraction yields stable algebraic updates.
  \item \textbf{Topological/Geometric Bases.} Choose $B_p$ as basis selectors; prime indices organize connectivity. Stable decompositions follow from Theorem~\ref{thm:convergence}.
  \item \textbf{Tensor Networks.} Let $B_p$ act on tensor cores/edges; the fixed point realizes a high-dimensional network with controlled growth.
  \item \textbf{Reinforcement Learning.} Take $T$ as weights, $F$ as TD/gradient drive, and $B_p$ as feature filters. Stability ensures convergent Q/policy updates.
  \item \textbf{Kernel-Level Scheduling.} Model dispatch channels as $B_p$ and resource demand as $F$ for a multiplicity-driven scheduler (Section~\ref{sec:kernel_multiplicity}).
\end{enumerate}

\begin{theorem}[Fractal Stability via Prime-Indexed IFS]
\label{thm:fractal_systems}
If $H=\prod_{j} H_j$ with componentwise contractions $\{T\mapsto F+\mathcal{K}T\}$ whose constants are $<1$, then the induced self-map on $H$ is a contraction; the unique attractor $T_\infty$ is the limit of a prime-indexed iterated function system (IFS). Hierarchical prime partitioning induces self-similar structure in the limit.
\end{theorem}

\subsection{Refinements: Compression, Hierarchies, Causality, and Performance}
\paragraph{Signature compression (checkpointing).}
At epoch $k$, store $c_k\!\gets\!e$ and reset $e\!\gets\!0$. The rolled signature $\overline{e}=\sum_{k} c_k + e$ equals the uncheckpointed signature, preserving audits exactly.

\paragraph{Exact/approximate/log-space conservation.}
Prefer exact exponent equality $\sum_i e(X_i)=e(Y)$. For noise, allow $\|\sum_i e(X_i)-e(Y)\|_\infty\le\varepsilon$. If only $M(e)$ is accessible, use $\sum_i\sum_p e_p(X_i)\log p \approx \sum_p e_p(Y)\log p$ with tolerance $\eta$.

\paragraph{Hierarchical rollups.}
For a module tree $\mathcal{T}$ with local signatures $e[u]$, define $E[u]=e[u]+\sum_{v\in\mathrm{Children}(u)}E[v]$. If primitives are conservative, then $E[u]$ matches $u$'s external interface; conservation at the root implies global conservation.

\paragraph{Causality (optional).}
Augment PETC with a DAG $G=(V,E)$; label each edge with its prime increment. The justification set for a node (state) is the multiset union along any source–node path; its signature equals the state’s PETC signature.

\subsection{ACE: Exact Weighted-\texorpdfstring{$\ell_1$}{l1} Projection}
We enforce $\sum_p b_p|w_p|\le\tau<1$ via the exact projection onto the weighted ball. Let $z_p=b_p|w_p|$. If $\sum_p z_p\le\tau$ return $w$; else find $\theta\ge0$ with $\sum_p \max(z_p-\theta,0)=\tau$ and set
\[
w_p^\star=\mathrm{sign}(w_p)\,\frac{\max(z_p-\theta,0)}{b_p}.
\]
This is the weighted analogue of soft-thresholding; $\theta$ can be found by sorting $\{z_p\}$ and a cumulative scan.

\subsection{Kernel-Level Multiplicity Scheduler}
\label{sec:kernel_multiplicity}
Associate each queue/class with a prime $p$; maintain PETC signatures per process and per core. Use ACE to allocate per-prime weights $w_p$ under $\sum b_p|w_p|\le\tau<1$; update dispatch decisions via $T_{t+1}=F+\mathcal{K}T_t$ where $T$ aggregates runnable-state indicators and $F$ encodes exogenous demand. PETC provides exact provenance for merges/splits; ACE maintains stability and responsiveness through budgeted adjustments.

\subsection{Algorithms and Code (Runnable Sketches)}

\paragraph{PETC prime ledger with checkpointing and exact conservation.}
\begin{verbatim}
class PrimeLedger:
    def __init__(self):
        self.sig = {}          # dict[int -> int], sparse exponents e_p
        self.checkpoints = []  # list[dict[int -> int]]

    def add_event(self, p: int, mult: int = 1, sign: int = +1):
        self.sig[p] = self.sig.get(p, 0) + sign * mult
        if self.sig[p] == 0: self.sig.pop(p)

    def checkpoint(self):
        self.checkpoints.append(self.sig.copy())
        self.sig.clear()

    def rolled_signature(self):
        roll = {}
        def add_map(dst, src, sgn=+1):
            for p, e in src.items():
                dst[p] = dst.get(p, 0) + sgn * e
                if dst[p] == 0: dst.pop(p)
        for ck in self.checkpoints: add_map(roll, ck, +1)
        add_map(roll, self.sig, +1)
        return roll

def conservation_ok_exact(inputs, out) -> bool:
    acc = {}
    def add(dst, src, sgn=+1):
        for p, e in src.items():
            dst[p] = dst.get(p, 0) + sgn * e
            if dst[p] == 0: dst.pop(p)
    for L in inputs: add(acc, L.rolled_signature(), +1)
    add(acc, out.rolled_signature(), -1)
    return len(acc) == 0
\end{verbatim}

\paragraph{Optional log-space surrogate (tolerant).}
\begin{verbatim}
import math

def log_signature(sig: dict[int,int]) -> float:
    return sum(e * math.log(p) for p, e in sig.items())

def conservation_ok_log(inputs, out, eta: float = 1e-9):
    lhs = sum(log_signature(L.rolled_signature()) for L in inputs)
    rhs = log_signature(out.rolled_signature())
    return abs(lhs - rhs) <= eta
\end{verbatim}

\paragraph{ACE exact weighted-\(\ell_1\) projection and prime-indexed operator.}
\begin{verbatim}
import math
import numpy as np

def project_weighted_l1(w: dict[int,float], b: dict[int,float], tau: float):
    z = {p: b[p] * abs(w[p]) for p in w}
    S = sum(z.values())
    if S <= tau: return w
    items = sorted(z.items(), key=lambda kv: kv[1], reverse=True)
    cumsum, theta = 0.0, 0.0
    for k, (_, zk) in enumerate(items, start=1):
        cumsum += zk
        theta = (cumsum - tau) / k
        if k == len(items) or items[k][1] <= theta < items[k-1][1]:
            break
    w_star = {}
    for p in w:
        val = max(z[p] - theta, 0.0) / max(b[p], 1e-12)
        w_star[p] = math.copysign(val, w[p])
    return w_star

class PrimeAffine:
    def __init__(self, B_ops: dict[int, callable], b_norms: dict[int,float], tau: float = 0.9):
        self.B = B_ops     # p -> H->H
        self.b = b_norms   # p -> ||B_p||
        self.tau = tau
        self.w  = {p: 0.0 for p in B_ops}

    def set_weights(self, Lambda: dict[int,float], alpha: float):
        w = {p: Lambda.get(p, 0.0) * (p ** alpha) for p in self.B}
        self.w = project_weighted_l1(w, self.b, self.tau)

    def K_apply(self, T):
        acc = 0 * T
        for p, Bp in self.B.items():
            acc = acc + self.w[p] * Bp(T)
        return acc

    def step(self, T, F):
        return F + self.K_apply(T)
\end{verbatim}

\paragraph{End-to-end loop (with PETC logging and ACE budget check).}
\begin{verbatim}
def run_loop(T0, F, prime_affine: PrimeAffine, ledger: PrimeLedger,
             prime_of_update: int, steps: int = 100):
    T = T0
    for t in range(steps):
        ledger.add_event(prime_of_update, mult=1, sign=+1)   # PETC log
        T_next = prime_affine.step(T, F)                     # ACE step
        budget = sum(abs(prime_affine.w[p]) * prime_affine.b[p] for p in prime_affine.w)
        assert budget < 1.0, "ACE budget violated: not a contraction"
        T = T_next
    return T
\end{verbatim}

\subsection{Operational Checklist}
\begin{enumerate}
  \item \textbf{Ledger (PETC).} Assign primes to components/events; maintain signatures additively; enforce conservation. Use checkpointing and hierarchical rollups for scale; optionally attach a causality DAG.
  \item \textbf{Weights (ACE).} Set $w_p=\Lambda_p p^\alpha$ with $\alpha< -1$ (sufficient) and project onto $\sum b_p|w_p|\le\tau<1$ via weighted-$\ell_1$ projection.
  \item \textbf{Certificates.} Monitor budget $\sum b_p|w_p|$, and (optionally) gap/slope margins during learning.
  \item \textbf{Convergence.} By Theorems~\ref{thm:convergence}–\ref{thm:uniqueness}, $T_{t+1}=F+\mathcal{K}T_t$ (or its projected form) converges uniquely to a fixed point at rate $\|\mathcal{K}\|$.
\end{enumerate}

\appendix
\section{Handy Design Bound for $|k|$}\label{app:design-bound}
For uniform $\Lambda_p=\Lambda$ and $\alpha<-1$,
\begin{equation*}
  |k|=\Lambda \sum_{i=1}^{N}p_i^{\alpha}.
\end{equation*}
Solve $|k|=1-\varepsilon$ for $\Lambda$ to initialize with margin $\varepsilon$.

\paragraph{Reproducibility Checklist (suggested).} Provide code/configs; fix seeds; publish CSV logs for $|k|$, budget usage, and certificate traces; include README with environment and commit hash.

\section{Prime-Recursive Mechanical Stabilization: PETC \texorpdfstring{$\times$}{x} ACE Formalization}
\label{sec:pirtm_formal}

\paragraph{Attribution.}
\textsc{ACE} (Arithmetic Control Engine): Tyler Van Osdol. \;
\textsc{PETC} (Prime-Encoded Tensor Calculus): Ryan Van Gelder.

\subsection{Prime-Signature Ledger (PETC) for Unambiguous Retrospection}
Let $\mathbb{P}$ denote the set of primes. A \emph{prime signature} is a finitely supported vector
$e \in \mathbb{Z}^{(\mathbb{P})}$; we write $e=(e_p)_{p \in \mathbb{P}}$ with $e_p=0$ for all but finitely many $p$.
Define the \emph{multiplicity map}
\[
  M(e) \;=\; \prod_{p\in\mathbb{P}} p^{e_p} \;\in\; \mathbb{Q}_{>0}.
\]
Operations are componentwise: composition of events is $e+e'$, undo/consumption is $e-e'$, dualization is $-e$.
This realizes the slogan ``\emph{multiply primes now, factor later}'' as \emph{add exponents now, read them later}:
unique factorization ensures that $e$ (hence the full interaction multiset) is recoverable without ambiguity.

\begin{definition}[Conservation Certificate]
For any typed operation $\Phi$ with inputs $X_1,\dots,X_r$ and output $Y$, the PETC \emph{conservation check} requires
\(
  e(Y) \;=\; \sum_{i=1}^r e(X_i).
\)
A system is \emph{ledger-consistent} if the equality holds for all executed $\Phi$.
\end{definition}

\begin{proposition}[Lossless Retrospection]
If a system is ledger-consistent, then for every state the accumulated signature $e$ determines the exact
multiset of prime-tagged interactions that occurred (with multiplicities), via the valuation $e_p = v_p(M(e))$.
\end{proposition}

\noindent\textbf{Complexity.}
Maintain $e$ as a sparse map (prime $\mapsto$ exponent). Updates and audits are $O(|\mathrm{supp}(e)|)$; no integer factorization is required.

\subsection{Prime-Indexed Affine Recursion (ACE-Certified)}
Let $H$ be a Banach space of tensors (axes carry PETC types).
For each prime $p$ choose a bounded linear ``channel'' $B_p:H\!\to\!H$ with $\|B_p\|\le b_p$.
Choose per-prime weights $w_p=\Lambda_p\,p^{\alpha}$ (design knob: $\Lambda_p\in\mathbb{R}$, typically $\Lambda_p\ge 0$, and $\alpha< -1$ for decay).
Define the gain operator
\[
  \mathcal{K}\;=\;\sum_{p\in P_N} w_p\,B_p,
\qquad
  \|\mathcal{K}\| \;\le\; \sum_{p\in P_N} |w_p|\,b_p.
\]
The \emph{meta-recursion} is the affine map
\begin{equation}
  \label{eq:affine}
  T_{t+1} \;=\; F \;+\; \mathcal{K}\,T_t,
\end{equation}
which collapses your original sum into an operator form while keeping primes as explicit channels.

\begin{assumption}[ACE Contraction Budget]\label{as:budget}
$\displaystyle \sum_{p\in P_N} |w_p|\,b_p \;<\;1.$
\end{assumption}

\begin{theorem}[Fixed Point, Uniqueness, and Linear Rate]
Under Assumption~\ref{as:budget}, the recursion \eqref{eq:affine} is a strict contraction on $(H,\|\cdot\|)$.
Hence there exists a unique fixed point $T_\infty = (I-\mathcal{K})^{-1}F$ and for any $T_0$
\[
  \|T_{t}-T_\infty\| \;\le\; \|\mathcal{K}\|^{\,t}\,\|T_0-T_\infty\|.
\]
In the scalar special case $T\in\mathbb{R}$ and $\mathcal{K}=k\in(-1,1)$, $T_\infty=\frac{F}{1-k}$.
\end{theorem}

\begin{remark}[Safe Design Regime]
A sufficient (not necessary) recipe is $\alpha< -1$, $\Lambda_p\ge 0$, and a global cap $\Lambda_p\le \bar\Lambda$ chosen so that
$\sum_{p\in P_N} \Lambda_p p^{\alpha} b_p < 1$. Increase $N$ or tighten $\bar\Lambda$ to keep the margin.
\end{remark}

\paragraph{ACE Safety Certificates.}
Let $\delta_S>0$ be a certified lower bound on the system's stability gap in the absence of $\mathcal{K}$,
and let $L_p$ bound the local schedule slope contributed by channel $B_p$.
Define
\[
  \mathrm{GapLB}(w) := \delta_S - 2\sum_{p\in P_N} b_p |w_p|, 
  \qquad
  \mathrm{SlopeUB}(w) := \sum_{p\in P_N} L_p |w_p|.
\]
Operate under the constraints $\mathrm{GapLB}(w)>0$ and $\mathrm{SlopeUB}(w)<1$.
These provide online, quantitative margins for stability and non-expansiveness during learning.

\subsection{Algorithms and Code Snippets}

\paragraph{PETC Prime Ledger (lossless history).}
\begin{verbatim}
class PrimeLedger:
    def __init__(self):
        self.sig = {}  # dict[int -> int], sparse exponents e_p

    def add_event(self, p: int, mult: int = 1, sign: int = +1):
        # sign=+1 for occurrence, -1 for consumption/undo
        self.sig[p] = self.sig.get(p, 0) + sign * mult
        if self.sig[p] == 0:
            self.sig.pop(p)

    def valuation(self, p: int) -> int:
        return self.sig.get(p, 0)

    def multiplicity_Q(self):
        # returns numerator, denominator of product p^{e_p}
        num, den = 1, 1
        for p, e in self.sig.items():
            if e >= 0: num *= p ** e
            else:      den *= p ** (-e)
        return num, den

def conservation_ok(inputs, output) -> bool:
    # Check PETC conservation: sum(e_in) == e_out
    acc = {}
    def add_into(dst, src, sgn=+1):
        for p, e in src.sig.items():
            dst[p] = dst.get(p, 0) + sgn * e
            if dst[p] == 0: dst.pop(p)
    for L in inputs: add_into(acc, L, +1)
    add_into(acc, output, -1)
    return len(acc) == 0
\end{verbatim}

\paragraph{ACE-Certified Affine Update with Budget Projection.}
We implement a simple (conservative) weighted-$\ell_1$ shrinkage to enforce the budget
$\sum_p b_p |w_p| \le \tau < 1$. Replace \texttt{shrink\_wl1} with an exact weighted-$\ell_1$ projection if needed.

\begin{verbatim}
import numpy as np

def shrink_wl1(w, b, tau):
    # Conservative shrink: scale if the budget is exceeded
    # w, b are dict[int->float]. Returns dict[int->float].
    budget = sum(abs(w[p]) * b[p] for p in w)
    if budget <= tau: return w
    scale = tau / budget
    return {p: scale * w[p] for p in w}

class PrimeAffine:
    def __init__(self, B_ops, b_norms, tau=0.9):
        # B_ops: dict[int -> callable H->H], b_norms: dict[int->float] with ||B_p||<=b_p
        self.B = B_ops
        self.b = b_norms
        self.tau = tau
        self.w  = {p: 0.0 for p in B_ops}  # start neutral

    def set_weights(self, Lambda, alpha):
        # Lambda: dict[int->float], alpha< -1 recommended
        w = {p: Lambda.get(p, 0.0) * (p ** alpha) for p in self.B}
        self.w = shrink_wl1(w, self.b, self.tau)  # enforce ACE budget

    def K_apply(self, T):
        acc = 0 * T
        for p, Bp in self.B.items():
            acc = acc + self.w[p] * Bp(T)
        return acc

    def step(self, T, F):
        # One affine contraction step: T_{t+1} = F + K T_t
        return F + self.K_apply(T)
\end{verbatim}

\paragraph{End-to-End Loop with Ledger and Stability Checks.}
\begin{verbatim}
def run_loop(T0, F, prime_affine, ledger, prime_of_update, steps=100):
    T = T0
    for t in range(steps):
        # Log the update event in PETC:
        ledger.add_event(prime_of_update, mult=1, sign=+1)

        # ACE contraction step:
        T_next = prime_affine.step(T, F)

        # Optional: online margin monitoring
        budget = sum(abs(prime_affine.w[p]) * prime_affine.b[p] for p in prime_affine.w)
        assert budget < 1.0, "ACE budget violated: not a contraction"

        T = T_next
    return T
\end{verbatim}

\subsection{Design Pattern and Examples}
\begin{itemize}
  \item \textbf{RL / Control.} Let $T$ be a weight tensor (e.g., linear Q or policy params), $F$ the external drive
        (TD target / gradient), and $B_p$ prime-typed feature filters. PETC logs every gradient/apply step; ACE ensures
        $\sum b_p |w_p| < 1$ for stable convergence to $T_\infty$.
  \item \textbf{Tensor Networks / Algebraic Systems.} Choose $B_p$ as structured projectors or sparsity masks whose norms
        are known. PETC ensures typed contractions respect multiplicity conservation; ACE budgets keep the build-up stable.
  \item \textbf{Auditing \& Replay.} Given a terminal signature $e$, valuations $e_p$ give exact interaction counts
        (multiset multiplicities). The system can replay or attribute decisions without ambiguity.
\end{itemize}

\subsection{Takeaways (Operational Rules)}
\begin{enumerate}
  \item \textbf{Ledger:} assign each component/event a prime; update signatures additively; enforce PETC conservation.
  \item \textbf{Weights:} set $w_p=\Lambda_p p^\alpha$ with $\alpha< -1$ (sufficient) and cap $\Lambda_p$ to satisfy
        $\sum |w_p| b_p < 1$.
  \item \textbf{Certificates:} monitor $\mathrm{GapLB}(w)>0$ and $\mathrm{SlopeUB}(w)<1$ online.
  \item \textbf{Convergence:} $T_{t+1}=F+\mathcal{K}T_t$ contracts to $T_\infty=(I-\mathcal{K})^{-1}F$ with linear rate $\|\mathcal{K}\|$.
\end{enumerate}


\nocite{*}
\bibliographystyle{plain}
\bibliography{references}

\end{document}
